Rice leaf diseases require timely visual diagnosis because they can affect crop productivity and disease management decisions. Recent convolutional neural networks (CNNs) and transfer learning models have achieved strong performance in plant disease recognition; however, RGB-based classification may still be influenced by background conditions, illumination variability, or dataset-specific visual patterns. This paper presents an applied, quantitative, experimental, and descriptive-comparative study that evaluates the contribution of classical segmentation and explicit texture descriptors in lightweight CNN-based rice leaf disease classification. Four configurations were compared: a MobileNetV2 baseline trained on original RGB images, a CNN trained on segmented images, a CNN fused with explicit texture features, and a hybrid configuration combining segmentation and texture features. The dataset split was controlled using perceptual-hash grouping to avoid visually equivalent samples across training, validation, and test sets. Experiments were repeated across five random seeds and evaluated using accuracy, macro-precision, macro-recall, macro-F1, inference time, model size, and parameter count. The texture-fusion configuration obtained the highest mean macro-F1 score of 0.9993, followed by the segmentation-only model with 0.9976, the proposed hybrid model with 0.9976, and the baseline with 0.9969. Although the enhanced configurations achieved higher average scores than the baseline, the differences were not statistically significant across five seeds. These results suggest that texture and segmentation are useful components for controlled ablation, but the dataset presents a performance saturation scenario in which several lightweight configurations already achieve very high classification performance.